\section{Give each checksum its own byte domain}
PNG's chunk CRC covers the four type bytes followed by the payload. It excludes the
length and the stored CRC itself. Reuse the reflected CRC-32 reasoning from the
checksum chapter, including the initial and final complements. Compare the resulting
32-bit pattern with the four stored big-endian bytes. Skipped ancillary data must
still contribute to this calculation.

The zlib stream has a different checksum, Adler-32, over the \emph{inflated} bytes.
For PNG, those bytes still include the row filter selectors and residuals. They are
not the final RGB values. A chunk CRC can pass while Adler-32 fails: recomputing a
chunk CRC after corrupting the compressed payload does not repair the inner stream.
Conversely, a correct zlib stream does not excuse a damaged chunk envelope.

CRC and Adler checks detect many accidental corruptions. Neither establishes who
created a file. For this program their purpose is to reject inconsistent input before
publishing pixels. Private scratch buffers may already have changed when a late
checksum fails; framebuffer publication must wait until validation is complete.
